\ifdefined\gkver\else\def\gkver{student}\fi
\documentclass[\gkver]{gaokaozhenti}
\begin{document}

\chapter{2025年黑吉辽内}

\begin{enumerate}
\item
%% number: 1
%% paperName: 2025年黑吉辽内高考第 1 题 4 分
%% typeId: 1
%% type: 单选题
%% chapter: 直线运动
%% point: 路程与位移
%% method: 无
%% score: 4
%% degree: 850
%% duplicateId: 0
%% body:
书法课上，某同学临摹“力”字时，笔尖的轨迹如图中带箭头的实线所示。笔尖由 a 点经 b 点回到 a 点，则\xzanswer{A}

\onepicture[w=2.60cm]{figs/01.png}

\fourchoices[answer=A]
{该过程位移为 0}
{该过程路程为 0}
{两次过 a 点时速度方向相同}
{两次过 a 点时摩擦力方向相同}

%% answer: A
%% memo:
\memoanswer{
A．笔尖由 a 点经 b 点回到 a 点过程，初位置和末位置相同，位移为零，故 A 正确；

B．笔尖由 a 点经 b 点回到 a 点过程，轨迹长度不为零，则路程不为零，故 B 错误；

C．两次过 a 点时轨迹的切线方向不同，则速度方向不同，故 C 错误；

D．摩擦力方向与笔尖的速度方向相反，则两次过 a 点时摩擦力方向不同，故 D 错误。

故选 A。
}

\item
%% number: 2
%% paperName: 2025年黑吉辽内高考第 2 题 4 分
%% typeId: 1
%% type: 单选题
%% chapter: 气体性质
%% point: 查理定律
%% method: 无
%% score: 4
%% degree: 750
%% duplicateId: 0
%% body:
某同学冬季乘火车旅行，在寒冷的站台上从气密性良好的糖果瓶中取出糖果后拧紧瓶盖，将糖果瓶带入温暖的车厢内一段时间后，与刚进入车厢时相比，瓶内气体\xzanswer{B}

\fourchoices[answer=B]
{内能变小}
{压强变大}
{分子的数密度变大}
{每个分子动能都变大}

%% answer: B
%% memo:
\memoanswer{
A．将糖果瓶带入温暖的车厢内一段时间后，温度升高，而理想气体内能只与温度相关，则内能变大，故 A 错误；

B．将糖果瓶带入温暖的车厢过程，气体做等容变化，根据查理定律可知温度升高，压强变大，故 B 正确；

C．气体分子数量不变，气体体积不变，则分子的数密度不变，故 C 错误；

D．温度升高，气体分子的平均动能增大，但不是每个分子的动能都增大，故 D 错误。

故选 B。
}

\item
%% number: 3
%% paperName: 2025年黑吉辽内高考第 3 题 4 分
%% typeId: 1
%% type: 单选题
%% chapter: 光学
%% point: 全反射
%% method: 无
%% score: 4
%% degree: 750
%% duplicateId: 0
%% body:
如图，利用液导激光技术加工器件时，激光在液束流与气体界面发生全反射。若分别用甲、乙两种液体形成液束流，甲的折射率比乙的大，则\xzanswer{D}

\onepicture[w=4.80cm]{figs/03.png}

\fourchoices[answer=D]
{激光在甲中的频率大}
{激光在乙中的频率大}
{用甲时全反射临界角大}
{用乙时全反射临界角大}

%% answer: D
%% memo:
\memoanswer{
AB．激光在不同介质中传播时，其频率不变，故 AB 错误；

CD．根据 $\sin C = \frac{1}{n}$，甲的折射率比乙的大，则用乙时全反射临界角大，故 C 错误，D 正确。

故选 D。
}

\item
%% number: 4
%% paperName: 2025年黑吉辽内高考第 4 题 4 分
%% typeId: 1
%% type: 单选题
%% chapter: 电场
%% point: 电容
%% method: 无
%% score: 4
%% degree: 550
%% duplicateId: 0
%% body:
如图，某压力传感器中平行板电容器内的绝缘弹性结构是模仿犰狳设计的，逐渐增大施加于两极板压力 $F$ 的过程中，$F$ 较小时弹性结构易被压缩，极板间距 $d$ 容易减小；$F$ 较大时弹性结构闭合，$d$ 难以减小。将该电容器充电后断开电源，极板间电势差 $U$ 与 $F$ 的关系曲线可能正确的是\xzanswer{D}

\onepicture[w=6.80cm]{figs/04a.png}

\fourchoices[answer=D, ispicture=true, h=2.40cm]
{figs/04b.png}
{figs/04c.png}
{figs/04d.png}
{figs/04e.png}

%% answer: D
%% memo:
\memoanswer{
根据公式 $Q = CU$ 和电容的决定式 $C = \frac{\varepsilon S}{4\pi kd}$ 可得 $U = \frac{4\pi kQ}{\varepsilon S}d$。

根据题意 $F$ 较小时易被压缩，故可知当 $F$ 较小时，随着 $F$ 的增大，$d$ 在减小，且减小的越来越慢，与电源断开后 $Q$ 不变，故此时极板间的电势差 $U$ 在减小，且减小的越来越慢；当 $F$ 增大到一定程度时，再增大 $F$ 后，$d$ 基本不变，故此时 $U$ 保持不变，结合图像，最符合情境的是 D 选项。

故选 D。
}

\item
%% number: 5
%% paperName: 2025年黑吉辽内高考第 5 题 4 分
%% typeId: 1
%% type: 单选题
%% chapter: 机械波
%% point: 干涉衍射
%% method: 无
%% score: 4
%% degree: 550
%% duplicateId: 0
%% body:
平衡位置在同一水平面上的两个振动完全相同的点波源，在均匀介质中产生两列波。若波峰用实线表示，波谷用虚线表示，P 点位于其最大正位移处，曲线 ab 上的所有点均为振动减弱点，则下列图中可能满足以上描述的是\xzanswer{C}

\fourchoices[answer=C, ispicture=true, h=2.60cm]
{figs/05a.png}
{figs/05b.png}
{figs/05c.png}
{figs/05d.png}

%% answer: C
%% memo:
\memoanswer{
根据题意 P 点位于其最大正位移处，故可知此时 P 点位于两列波的波峰与波峰相交处，BD 错误；根据干涉规律可知，相邻波峰与波峰，波谷与波谷连线上的点都是加强点，故 A 图像中的曲线 ab 上的点存在振动加强点，不符合题意。

故选 C。
}

\item
%% number: 6
%% paperName: 2025年黑吉辽内高考第 6 题 4 分
%% typeId: 1
%% type: 单选题
%% chapter: 直线运动
%% point: 运动的分解
%% method: 无
%% score: 4
%% degree: 450
%% duplicateId: 0
%% body:
如图，趣味运动会的“聚力建高塔”活动中，两长度相等的细绳一端系在同一塔块上，两名同学分别握住绳的另一端，保持手在同一水平面以相同速率 $v$ 相向运动。为使塔块沿竖直方向匀速下落，则 $v$\xzanswer{B}

\onepicture[w=6.00cm]{figs/06.png}

\fourchoices[answer=B]
{一直减小}
{一直增大}
{先减小后增大}
{先增大后减小}

%% answer: B
%% memo:
\memoanswer{
设两边绳与竖直方向的夹角为 $\theta$，塔块沿竖直方向匀速下落的速度为 $v_{\text{块}}$，将 $v_{\text{块}}$ 沿绳方向和垂直绳方向分解，将 $v$ 沿绳子方向和垂直绳方向分解，可得 $v_{\text{块}}\cos\theta = v\sin\theta$，解得 $v = \frac{v_{\text{块}}}{\tan\theta}$。由于塔块匀速下落时 $\theta$ 在减小，故可知 $v$ 一直增大。

故选 B。
}

\item
%% number: 7
%% paperName: 2025年黑吉辽内高考第 7 题 4 分
%% typeId: 1
%% type: 单选题
%% chapter: 功和能
%% point: 动能定理
%% method: 无
%% score: 4
%% degree: 450
%% duplicateId: 0
%% body:
如图，光滑绝缘水平面 AB 与竖直面内光滑绝缘半圆形轨道 BC 在 B 点相切，轨道半径为 $r$，圆心为 O，O、A 间距离为 $3r$。原长为 $2r$ 的轻质绝缘弹簧一端固定于 O 点，另一端连接一带正电的物块。空间存在水平向右的匀强电场，物块所受的电场力与重力大小相等。物块在 A 点左侧释放后，依次经过 A、B、C 三点时的动能分别为 $E_{kA}$、$E_{kB}$、$E_{kC}$，则\xzanswer{C}

\onepicture[w=5.60cm]{figs/07.png}

\fourchoices[answer=C]
{$E_{kA} < E_{kB} < E_{kC}$}
{$E_{kB} < E_{kA} < E_{kC}$}
{$E_{kA} < E_{kC} < E_{kB}$}
{$E_{kC} < E_{kA} < E_{kB}$}

%% answer: C
%% memo:
\memoanswer{
由题意可得 A 点弹簧伸长量为 $r$，B 点和 C 点弹簧压缩量为 $r$，即三个位置弹簧弹性势能相等，则由 A 到 B 过程中弹簧弹力做功为零，电场力做正功，动能增加，$E_{kB} > E_{kA}$；同理 B 到 C 过程中弹簧弹力和电场力做功都为零，重力做负功，则动能减小，$E_{kB} > E_{kC}$；由 A 到 C 全过程则有

$qEl_{AB} - mgl_{BC} = E_{kC} - E_{kA} > 0$

因此 $E_{kB} > E_{kC} > E_{kA}$

故选 C。
}

\item
%% number: 8
%% paperName: 2025年黑吉辽内高考第 8 题 6 分
%% typeId: 2
%% type: 多选题
%% chapter: 原子和原子核
%% point: 天然放射性
%% method: 无
%% score: 6
%% degree: 650
%% duplicateId: 0
%% body:
某理论研究认为，\ce{^{100}_{42}Mo} 原子核可能发生双 $\beta$ 衰变，衰变方程为 \ce{^{100}_{42}Mo} $\to$ \ce{^{A}_{44}Ru} + $y$ \ce{^{0}_{-1}e}。处于第二激发态的 \ce{^{A}_{44}Ru} 原子核先后辐射能量分别为 $0.590 8 \UMeV$ 和 $0.539 5 \UMeV$ 的 $\gamma_{1}$、$\gamma_{2}$ 两光子后回到基态。下列说法正确的是\xzanswer{ABC}

\fourchoices[answer=ABC]
{$A = 100$}
{$y = 2$}
{$\gamma_{1}$ 的频率比 $\gamma_{2}$ 的大}
{$\gamma_{1}$ 的波长比 $\gamma_{2}$ 的大}

%% answer: ABC
%% memo:
\memoanswer{
AB．由核反应方程质量数和电荷数守恒可得 $100 = A + 0$，$42 = 44 - y$，解得 $A = 100$，$y = 2$，AB 正确；

CD．由题可得 $\gamma_{1}$ 光子的能量大于 $\gamma_{2}$ 光子的能量，光子的能量公式 $E = h\nu$，波长 $\lambda = \frac{c}{\nu}$，可得 $\gamma_{1}$ 的频率大于 $\gamma_{2}$ 的频率，$\gamma_{1}$ 的波长小于 $\gamma_{2}$ 的波长，C 正确，D 错误；

故选 ABC。
}

\item
%% number: 9
%% paperName: 2025年黑吉辽内高考第 9 题 6 分
%% typeId: 2
%% type: 多选题
%% chapter: 电磁感应
%% point: 切割磁感线电动势
%% method: 无
%% score: 6
%% degree: 550
%% duplicateId: 0
%% body:
如图，“$\llcorner$”形导线框置于磁感应强度大小为 $B$、水平向右的匀强磁场中。线框相邻两边均互相垂直，各边长均为 $l$。线框绕 b、e 所在直线以角速度 $\omega$ 顺时针匀速转动，be 与磁场方向垂直。$t = 0$ 时，abef 与水平面平行，则\xzanswer{AB}

\onepicture[w=3.90cm]{figs/09.png}

\fourchoices[answer=AB]
{$t = 0$ 时，电流方向为 abcdefa}
{$t = 0$ 时，感应电动势为 $Bl^{2}\omega$}
{$t = \frac{\pi}{\omega}$ 时，感应电动势为 0}
{$t = 0$ 到 $t = \frac{\pi}{\omega}$ 过程中，感应电动势平均值为 0}

%% answer: AB
%% memo:
\memoanswer{
AB．线框旋转切割磁场产生电动势的两条边为 cd 和 af，$t = 0$ 时刻 cd 边速度与磁场方向平行，不产生电动势，因此此时 af 边切割产生电动势，由右手定则可知电流方向为 abcdefa，电动势为 $E = Blv = Bl\omega l = Bl^{2}\omega$，AB 正确；

C．$t = \frac{\pi}{\omega}$ 时，线框旋转 $180^\circ$，此时依旧是 af 边切割磁场产生电动势，感应电动势不为零，C 错误；

D．$t = 0$ 到 $t = \frac{\pi}{\omega}$ 时，线框 abef 的磁通量变化量为零，线框 bcde 的磁通量变化量为 $\Delta\Phi = 2BS = 2Bl^{2}$

由法拉第电磁感应定律可得平均电动势为 $E = \frac{\Delta\Phi}{\Delta t} = \frac{2B\omega l^{2}}{\pi}$，D 错误。

故选 AB。
}

\item
%% number: 10
%% paperName: 2025年黑吉辽内高考第 10 题 6 分
%% typeId: 2
%% type: 多选题
%% chapter: 运动和力
%% point: 牛顿第二定律
%% method: 整体与隔离
%% score: 6
%% degree: 350
%% duplicateId: 0
%% body:
如图~\ref{2025黑吉辽内10a}，倾角为 $\theta$ 的足够长斜面放置在粗糙水平面上。质量相等的小物块甲、乙同时以初速度 $v_{0}$ 沿斜面下滑，甲、乙与斜面的动摩擦因数分别为 $\mu_{1}$、$\mu_{2}$，整个过程中斜面相对地面静止。甲和乙的位置 $x$ 与时间 $t$ 的关系曲线如图~\ref{2025黑吉辽内10b}所示，两条曲线均为抛物线，乙的 $x-t$ 曲线在 $t = t_{0}$ 时切线斜率为 0，则\xzanswer{AD}

\twopicture[labelA=2025黑吉辽内10a, labelB=2025黑吉辽内10b, wA=3.60cm, wB=3.90cm]
{figs/10a.png}
{figs/10b.png}

\fourchoices[answer=AD]
{$\mu_{1} + \mu_{2} = 2\tan\theta$}
{$t = t_{0}$ 时，甲的速度大小为 $3v_{0}$}
{$t = t_{0}$ 之前，地面对斜面的摩擦力方向向左}
{$t = t_{0}$ 之后，地面对斜面的摩擦力方向向左}

%% answer: AD
%% memo:
\memoanswer{
B．位置 $x$ 与时间 $t$ 的图像的斜率表示速度，甲乙两个物块的曲线均为抛物线，则甲物体做匀加速运动，乙物体做匀减速运动，在 $t_{0}$ 时间内甲乙的位移可得 $x_{\text{甲}} = \frac{v_{0} + v}{2}t_{0} = 3x_{0}$，$x_{\text{乙}} = \frac{v_{0} + 0}{2}t_{0} = x_{0}$

可得 $t_{0}$ 时刻甲物体的速度为 $v = 2v_{0}$，B 错误；

A．甲物体的加速度大小为 $a_{1} = \frac{v - v_{0}}{t_{0}}$

乙物体的加速度大小为 $a_{2} = \frac{v_{0}}{t_{0}}$

由牛顿第二定律可得甲物体 $mg\sin\theta - \mu_{1}mg\cos\theta = ma_{1}$

同理可得乙物体 $\mu_{2}mg\cos\theta - mg\sin\theta = ma_{2}$

联立可得 $\mu_{1} + \mu_{2} = 2\tan\theta$，A 正确；

C．设斜面的质量为 $M$，取水平向左为正方向，由系统牛顿第二定律可得 $f = ma_{1}\cos\theta - ma_{2}\cos\theta = 0$

则 $t = t_{0}$ 之前，地面和斜面之间摩擦力为零，C 错误；

D．$t = t_{0}$ 之后，乙物体保持静止，甲物体继续沿斜面向下加速，由系统牛顿第二定律可得 $f = ma_{1}\cos\theta$

即地面对斜面的摩擦力向左，D 正确。

故选 AD。
}

\item
%% number: 11
%% paperName: 2025年黑吉辽内高考第 11 题 8 分
%% typeId: 4
%% type: 实验题
%% chapter: 电路
%% point: 伏安法测电阻
%% method: 无
%% score: 8
%% degree: 550
%% duplicateId: 0
%% body:
在测量某非线性元件的伏安特性时，为研究电表内阻对测量结果的影响，某同学设计了如图~\ref{2025黑吉辽内11a}所示的电路。选择多用电表的直流电压挡测量电压。实验步骤如下：
\begin{steps}
\item
滑动变阻器滑片置于适当位置，闭合开关；
\item
表笔分别连 a、b 接点，调节滑片位置，记录电流表示数 $I$ 和 a、b 间电压 $U_{ab}$；
\item
表笔分别连 a、c 接点，调节滑片位置，使电流表示数仍为 $I$，记录 a、c 间电压 $U_{ac}$；
\item
表笔分别连 b、c 接点，调节滑片位置，使电流表示数仍为 $I$，记录 b、c 间电压 $U_{bc}$，计算 $U_{ac} - U_{bc}$；
\item
改变电流，重复步骤②③④，断开开关。
\end{steps}

作出 $I-U_{ab}$、$I-U_{ac}$ 及 $I-(U_{ac} - U_{bc})$ 曲线如图~\ref{2025黑吉辽内11b}所示。

回答下列问题：
\begin{enumerate}
\item
将多用电表的红、黑表笔插入正确的插孔，测量 a、b 间的电压时，红表笔应连\tkanswer{a}接点（填“a”或“b”）；
\item
若多用电表选择开关旋转到直流电压挡“$0.5 \UV$”位置，电表示数如图~\ref{2025黑吉辽内11c}所示，此时电表读数为\tkanswer[1.8cm]{0.376$\sim$0.378}$\UV$（结果保留三位小数）；
\item
图~\ref{2025黑吉辽内11b}中乙是\tkanswer[1.6cm]{$I-U_{ac}$}（填“$I-U_{ab}$”或“$I-U_{ac}$”）曲线；
\item
实验结果表明，当此元件阻值较小时，\tkanswer{甲}（填“甲”或“乙”）曲线与 $I-(U_{ac} - U_{bc})$ 曲线更接近。
\end{enumerate}

\threepicture[labelA=2025黑吉辽内11a, labelB=2025黑吉辽内11b, labelC=2025黑吉辽内11c, wA=4.20cm, wB=3.40cm, wC=4.60cm, align=right]
{figs/11a.png}
{figs/11b.png}
{figs/11c.png}

%% answer:
\jdanswer{
\begin{enumerate}
\item
a

\item
$0.376\sim0.378$

\item
$I-U_{ac}$

\item
甲
\end{enumerate}
}
%% memo:
\memoanswer{
（1）电流从红表笔流入，黑表笔流出，故测量 a、b 间的电压时，红表笔应连 a 接点。

（2）$0.5 \UV$ 的直流电压挡，分度值为 $0.01 \UV$，由图可知此时电压表读数为 $0.377 \UV$。

（3）由图可知，当表笔分别连 a、c 接点时测得是元件和电流表两端的电压和电流，则 $R_{ac} = \frac{U_{ac}}{I}$

当表笔分别连 a、b 接点时测得是元件两端的电压和电流，则 $R_{ab} = \frac{U_{ab}}{I}$

由于 $R_{ac} > R_{ab}$

所以相同电流情况下，$U_{ac} > U_{ab}$

故图（b）中乙是 $I-U_{ac}$ 曲线。

（4）由题意可知，$I-(U_{ac} - U_{bc})$ 图像测得是元件两端的电压和电流的关系，则实验结果表明，当此元件阻值较小时，甲曲线与 $I-(U_{ac} - U_{bc})$ 曲线更接近。
}

\item
%% number: 12
%% paperName: 2025年黑吉辽内高考第 12 题 8 分
%% typeId: 4
%% type: 实验题
%% chapter: 力物体的平衡
%% point: 三力平衡
%% method: 无
%% score: 8
%% degree: 450
%% duplicateId: 0
%% body:
某兴趣小组设计了一个可以测量质量的装置。如图~\ref{2025黑吉辽内12a}，细绳 1、2 和橡皮筋相连于一点，绳 1 上端固定在 A 点，绳 2 下端与水杯相连，橡皮筋的另一端与绳套相连。

为确定杯中物体质量 $m$ 与橡皮筋长度 $x$ 的关系，该小组逐次加入等质量的水，拉动绳套，使绳 1 每次与竖直方向夹角均为 $30^\circ$ 且橡皮筋与绳 1 垂直，待装置稳定后测量对应的橡皮筋长度。根据测得数据作出 $x-m$ 关系图线，如图~\ref{2025黑吉辽内12b}所示。

回答下列问题：
\begin{enumerate}
\item
将一芒果放入此空杯，按上述操作测得 $x = 11.60 \Ucm$，由图~\ref{2025黑吉辽内12b}可知，该芒果的质量 $m_{0} =$\tkanswer[1.4cm]{106}$\Ug$（结果保留到个位）。若杯中放入芒果后，绳 1 与竖直方向夹角为 $30^\circ$ 但与橡皮筋不垂直，由图像读出的芒果质量与 $m_{0}$ 相比\tkanswer{偏大}（填“偏大”或“偏小”）。
\item
另一组同学利用同样方法得到的 $x-m$ 图像在后半部分弯曲，下列原因可能的是\tkanswer{C}。

A．水杯质量过小

B．绳套长度过大

C．橡皮筋伸长量过大，弹力与其伸长量不成正比
\item
写出一条可以使上述装置测量质量范围增大的措施\tkanswer{减小细线与竖直方向的夹角}。
\end{enumerate}

\twopicture[labelA=2025黑吉辽内12a, labelB=2025黑吉辽内12b, wA=3.20cm, wB=7.40cm, align=right]
{figs/12a.png}
{figs/12b.png}

%% answer:
\jdanswer{
\begin{enumerate}
\item
106，偏大

\item
C

\item
减小细线与竖直方向的夹角
\end{enumerate}
}
%% memo:
\memoanswer{
（1）操作测得 $x = 11.60 \Ucm$，由题图的图像坐标可知，该芒果的质量为 $106 \Ug$；

若杯中放入芒果后，绳 1 与竖直方向夹角为 $30^\circ$ 但与橡皮筋不垂直，根据共点力平衡可知橡皮筋的拉力变大，导致橡皮筋的长度偏大，若仍然根据图像读出芒果的质量与 $m_{0}$ 相比偏大。

\onepicture[w=3.00cm]{figs/12_an.png}

（2）另一组同学利用同样方法得到的 $x-m$ 图像在后半部分弯曲，可能是所测物体的质量过大，导致橡皮筋所受的弹力过大超过了弹簧的弹性限度，从而使橡皮筋弹力与其伸长量不成正比。

故选 C。

（3）根据共点力平衡条件可知，当减小绳子与竖直方向的夹角时，相同的物体质量对应橡皮筋的拉力较小，故相同的橡皮筋，减小细线与竖直方向的夹角可增大质量测量范围。
}

\item
%% number: 13
%% paperName: 2025年黑吉辽内高考第 13 题 10 分
%% typeId: 6
%% type: 计算题
%% chapter: 曲线运动
%% point: 斜抛运动
%% method: 无
%% score: 10
%% degree: 550
%% duplicateId: 0
%% body:
如图，一雪块从倾角 $\theta = 37^\circ$ 的屋顶上的 O 点由静止开始下滑，滑到 A 点后离开屋顶。O、A 间距离 $x = 2.5 \Um$，A 点距地面的高度 $h = 1.95 \Um$，雪块与屋顶的动摩擦因数 $\mu = 0.125$。不计空气阻力，雪块质量不变，取 $\sin 37^\circ = 0.6$，重力加速度大小 $g = 10 \Umsq$。求：
\begin{enumerate}
\item
雪块从 A 点离开屋顶时的速度大小 $v_{0}$；
\item
雪块落地时的速度大小 $v_{1}$，及其速度方向与水平方向的夹角 $\alpha$。
\end{enumerate}

\onepicture[w=5.20cm, align=right]{figs/13.png}

%% answer:
\jdanswer{
\begin{enumerate}
\item
$5 \Ums$

\item
$8 \Ums$，$60^\circ$
\end{enumerate}
}
%% memo:
\memoanswer{
（1）雪块在屋顶上运动过程中，由动能定理

$mgx\sin\theta - \mu mg\cos\theta\cdot x = \frac{1}{2}mv_{0}^{2} - 0$

代入数据解得雪块到 A 点速度大小为 $v_{0} = 5 \Ums$

（2）雪块离开屋顶后，做斜下抛运动，由动能定理

$mgh = \frac{1}{2}mv_{1}^{2} - \frac{1}{2}mv_{0}^{2}$

代入数据解得雪块到地面速度大小 $v_{1} = 8 \Ums$

速度与水平方向夹角 $\alpha$，满足

$\cos\alpha = \frac{v_{0}\cos\theta}{v_{1}} = \frac{5 \times 0.8}{8} = \frac{1}{2}$

解得 $\alpha = 60^\circ$
}

\item
%% number: 14
%% paperName: 2025年黑吉辽内高考第 14 题 12 分
%% typeId: 6
%% type: 计算题
%% chapter: 电磁感应
%% point: 法拉第电磁感应定律
%% method: 无
%% score: 12
%% degree: 450
%% duplicateId: 0
%% body:
如图~\ref{2025黑吉辽内14a}，固定在光滑绝缘水平面上的单匝正方形导体框 abcd，置于始终竖直向下的匀强磁场中，ad 边与磁场边界平行，ab 边中点位于磁场边界。导体框的质量 $m = 1 \Ukg$，电阻 $R = 0.5 \UO$、边长 $L = 1 \Um$。磁感应强度 $B$ 随时间 $t$ 连续变化，$0\sim1 \Us$ 内 $B-t$ 图像如图~\ref{2025黑吉辽内14b}所示。导体框中的感应电流 $I$ 与时间 $t$ 关系图像如图~\ref{2025黑吉辽内14c}所示，其中 $0\sim1 \Us$ 内的图像未画出，规定顺时针方向为电流正方向。
\begin{enumerate}
\item
求 $t = 0.5 \Us$ 时 ad 边受到的安培力大小 $F$；
\item
画出图~\ref{2025黑吉辽内14b}中 $1\sim2 \Us$ 内 $B-t$ 图像（无需写出计算过程）；
\item
从 $t = 2 \Us$ 开始，磁场不再随时间变化。之后导体框解除固定，给导体框一个向右的初速度 $v_{0} = 0.1 \Ums$，求 ad 边离开磁场时的速度大小 $v_{1}$。
\end{enumerate}

\threepicture[labelA=2025黑吉辽内14a, labelB=2025黑吉辽内14b, labelC=2025黑吉辽内14c, wA=3.30cm, wB=3.50cm, wC=3.90cm, align=right]
{figs/14a.png}
{figs/14b.png}
{figs/14c.png}

%% answer:
\jdanswer{
\begin{enumerate}
\item
$0.015 \UN$

\item
\onepicture[w=4.00cm]{figs/14_an.png}

\item
$0.01 \Ums$
\end{enumerate}
}
%% memo:
\memoanswer{
（1）由法拉第电磁感应定律

$E_{1} = \frac{\Delta\Phi}{\Delta t} = \frac{\Delta B\cdot\frac{1}{2}L^{2}}{\Delta t} = \frac{0.2 - 0.1}{1 - 0} \times \frac{1}{2} \times 1^{2} \UV = 0.05 \UV$

由闭合电路欧姆定律可知，$0\sim1 \Us$ 内线框中的感应电流大小为 $I_{1} = \frac{E_{1}}{R} = 0.1 \UA$

由图（b）可知，$t = 0.5 \Us$ 时磁感应强度大小为 $B_{0.5} = 0.15 \UT$

所以此时导线框 ad 的安培力大小为 $F = B_{0.5}I_{1}L = 0.15 \times 0.1 \times 1 \UN = 0.015 \UN$

（2）$0\sim1 \Us$ 内线框内的感应电流大小为 $I_{1} = 0.1 \UA$，根据楞次定律及安培定则可知感应电流方向为顺时针，由图（c）可知 $1\sim2 \Us$ 内的感应电流大小为 $I_{2} = 0.2 \UA$，方向为逆时针。根据欧姆定律可知 $1\sim2 \Us$ 内的感应电动势大小为 $E_{2} = I_{2}R = 0.1 \UV$

由法拉第电磁感应定律

$E_{2} = \frac{\Delta\Phi}{\Delta t} = \frac{\Delta B\cdot\frac{1}{2}L^{2}}{\Delta t} = 0.1 \UV$

可知 $1\sim2 \Us$ 内磁感应强度的变化率为 $\frac{\Delta B}{\Delta t} = \frac{B_{2} - B_{1}}{\Delta t} = 0.2 \UTs$

解得 $t = 2 \Us$ 时磁感应强度大小为 $B_{2} = 0.3 \UT$

方向垂直于纸面向里，故 $1\sim2 \Us$ 的磁场随时间变化图像见答案图。

（3）由动量定理可知 $-B_{2}\overline{I}L\Delta t = mv_{1} - mv_{0}$

其中 $q = \overline{I}\Delta t = \frac{\overline{E}}{R}\Delta t = \frac{\Delta\Phi}{R} = \frac{\frac{1}{2}B_{2}L^{2}}{R}$

联立解得 ad 经过磁场边界的速度大小为 $v_{1} = 0.01 \Ums$
}

\item
%% number: 15
%% paperName: 2025年黑吉辽内高考第 15 题 16 分
%% typeId: 6
%% type: 计算题
%% chapter: 磁场
%% point: 带电粒子在复合场中的运动
%% method: 无
%% score: 16
%% degree: 350
%% duplicateId: 0
%% body:
如图，在 $xOy$ 平面第一、四象限内存在垂直平面向里的匀强磁场，磁感应强度大小为 $B$，一带正电的粒子从 M（0，$-y_{0}$）点射入磁场，速度方向与 $y$ 轴正方向夹角 $\theta = 30^\circ$，从 N（0，$y_{0}$）点射出磁场。已知粒子的电荷量为 $q$（$q > 0$），质量为 $m$，忽略粒子重力及磁场边缘效应。
\begin{enumerate}
\item
求粒子射入磁场的速度大小 $v_{1}$ 和在磁场中运动的时间 $t_{1}$。
\item
若在 $xOy$ 平面内某点固定一负点电荷，电荷量为 $48q$，粒子质量取 $m = \frac{B^{2}y_{0}^{3}}{k}$（$k$ 为静电力常量），粒子仍沿（1）中的轨迹从 M 点运动到 N 点，求射入磁场的速度大小 $v_{2}$。
\item
在（2）问条件下，粒子从 N 点射出磁场开始，经时间 $t_{2}$ 速度方向首次与 N 点速度方向相反，求 $t_{2}$（电荷量为 $Q$ 的点电荷产生的电场中，取无限远处的电势为 0 时，与该点电荷距离为 $r$ 处的电势 $\varphi = \frac{kQ}{r}$）。
\end{enumerate}

\onepicture[w=3.00cm, align=right]{figs/15.png}

%% answer:
\jdanswer{
\begin{enumerate}
\item
$\frac{2qBy_{0}}{m}$，$\frac{\pi m}{3qB}$

\item
$\frac{6kq}{By_{0}^{2}}$

\item
$\frac{2\sqrt{3}\pi By_{0}^{3}}{3kq}$
\end{enumerate}
}
%% memo:
\memoanswer{
（1）作出正电荷在磁场中运动的轨迹，如图所示。

\onepicture[w=4.60cm]{figs/15_an1.png}

由几何关系可知，正电荷在磁场中做匀速圆周运动的半径为 $r = \frac{y_{0}}{\sin\theta} = 2y_{0}$

由洛伦兹力提供向心力 $qv_{1}B = m\frac{v_{1}^{2}}{r}$

解得正电荷的入射速度大小为 $v_{1} = \frac{2qBy_{0}}{m}$

正电荷在磁场中运动的周期为 $T = \frac{2\pi r}{v_{1}} = \frac{2\pi m}{qB}$

所以正电荷从 M 运动到 N 的时间为 $t_{1} = \frac{2\theta}{2\pi}T = \frac{\pi m}{3qB}$

（2）由题意可知，在 $xOy$ 平面内的负电荷在圆心 $O$ 处，由牛顿第二定律可知 $qv_{2}B + k\frac{48q^{2}}{r^{2}} = m\frac{v_{2}^{2}}{r}$，其中 $m = \frac{B^{2}y_{0}^{3}}{k}$

解得 $v_{2} = \frac{6kq}{By_{0}^{2}}$ 或 $-\frac{4kq}{By_{0}^{2}}$（舍去）

（3）在（2）的条件下，由题意可知，粒子从 N 点离开，仅在点电荷的作用下运动，粒子所需要的向心力 $m\frac{v^{2}}{r}$ 大于点电荷提供的库仑力 $k\frac{48q^{2}}{r^{2}}$，因此粒子无法做匀速圆周运动，即正电荷从 $N$ 点离开磁场后绕负电荷做椭圆运动，如图所示

\onepicture[w=5.00cm]{figs/15_an2.png}

由能量守恒定律得 $\frac{1}{2}mv_{2}^{2} - q\frac{k48q}{r_{2}} = \frac{1}{2}mv_{3}^{2} - q\frac{k48q}{r_{3}}$

由开普勒第二定律可知 $v_{2}r_{2} = v_{3}r_{3}$

其中 $r_{2} = 2y_{0}$

联立解得 $r_{3} = 6y_{0}$

由牛顿第二定律 $k\frac{48q^{2}}{\left(\frac{r_{2} + r_{3}}{2}\right)^{2}} = m\left(\frac{r_{2} + r_{3}}{2}\right)\frac{4\pi^{2}}{T^{2}}$

解得 $T = \frac{4\sqrt{3}\pi By_{0}^{3}}{3kq}$

故正电荷从 N 点离开磁场后到首次速度变为与 N 点的射出速度相反的时间为

$t_{2} = \frac{2\sqrt{3}\pi By_{0}^{3}}{3kq}$
}

\end{enumerate}

\end{document}
